\section{Global Notation}
\label{sec:global-notation}

The notation used throughout the proof is summarized here.  Every local
quantity introduced later is defined at its first use.

\small
\renewcommand{\arraystretch}{1.08}
\begin{longtable}{@{}p{.26\textwidth}p{.67\textwidth}@{}}
\toprule
symbol&meaning\\
\midrule
$H,O,V_i,e_{i,i+1},r_i,M_i$&
the unit hexagon, its center and vertices, a boundary edge, a radial segment,
and a radial midpoint\\
$S$&$\partial H\cup\bigcup_i r_i$, the full one-dimensional skeleton\\
$U_C,U_i$&the open center triangle and the open vertex triangle containing
$V_i$\\
$T_C,T_i$&the closures of $U_C,U_i$\\
$\CEzero,\CEone,\CEtwo$&the three center-triangle types, according to the
number of positive-length boundary traces\\
$\Vdzero,\Vdone,\Vdtwo,\Tlike$&the four vertex-triangle types determined by
the pair $(o,n)$ in Section~\ref{sec:introduction}\\
$A_i,B_i,C_i$&the maximal reaches of $T_i$ toward $V_{i-1},V_{i+1},O$\\
$(a_i,b_i,c_i)$&selected lower bounds for $(A_i,B_i,C_i)$\\
$N_+$&the number of indices for which $A_i+B_i>1$\\
$\Nsp$&the number of vertex triangles of type $\Vdone,\Vdtwo,$ or $\Tlike$\\
$\Ngap$&the number of positive center traces containing a boundary gap\\
$\mathcal A$&the admissible set of triples of reach lower bounds\\
$M_c(a)$&the maximum possible forward reach when the backward and radial
reaches are at least $a$ and $c$\\
$\overline M_c(a)$&the corresponding forward-reach cap for a
nonsupercritical vertex triangle\\
$\Phi_c(a)=1-\overline M_c(a)$&the lower bound propagated to the next
backward reach on a center-free edge\\
$\Lambda(K)$&the least side length of a closed equilateral triangle containing
$K$\\
\bottomrule
\end{longtable}
\normalsize
